\section{行业价格 $\beta$ 已得到交叉验证，但不是公司 ASP}

基于 L3 行业报价的周期背景，电池级碳酸锂在 2026H1 的价格环境明显强于 2025 年末。已归档公开报价显示，2025 年 12 月 31 日为 11.85 万元/吨，2026 年 1 月和 2 月的月度背景值为 15.60 万元/吨和 14.96 万元/吨；2026 年 7 月 1 日可捕获的现货区间为 15.51--16.51 万元/吨，均价 16.01 万元/吨。两个单点的机械比较为 +35.1\%，但它不是 H1 均价涨幅，更不是雅化实际 ASP 或毛利率。

\begin{exhibitbox}[锂价：可用作事件背景，不能直接计入雅化损益]
\begin{tabularx}{\exhibitboxwidth}{L{2.4cm}L{3.6cm}L{1.5cm}X}
\toprule
时间 & 电池级碳酸锂公开信息 & 资料等级 & 正确用途 \\
\midrule
2025年12月31日 & 11.85 万元/吨 & L3 & 周期起点参考 \\
2026年1月 / 2月 & 15.60 / 14.96 万元/吨月均 & L3 & Q1 价格环境背景 \\
2026年3--5月 & 月均环比约 +5\% / +6\% / +12\% & L3 & 说明价格弹性与非线性风险 \\
2026年6月 & 中枢约 16 万元/吨附近 & L3 & Q2 末方向性锚 \\
2026年7月1日 & 15.51--16.51 万元/吨，均价 16.01 万元/吨 & L3 & 本次抓取最新明确点，非 7 月 23 日实时价 \\
\bottomrule
\end{tabularx}
\end{exhibitbox}
\sourcenote{SMM 2026H1 锂市场回顾（2026-07-10）；中国有色网转载《2026年7月1日上海金属现货报价（钴锂）》（2026-07-01）；SMM 2025 年末现货回顾。所有价格均为行业资料，非雅化结算价。}

终端需求没有消失：人民网转述中汽协数据称 2026H1 新能源车销量为 744.6 万辆。但供给、进口和库存使价格没有单边逻辑：SMM 回顾称国内碳酸锂 H1 产出约 62.2 万吨，5 月累计进口 15.3 万吨、同比 +53\%，6 月仓单约 5 万吨。对雅化而言，这些数字的意义是 H2 的价格情景仍有上下行，而不是把新能源汽车销量换算为公司订单。

\section{同业 Q1：确认行业回暖，不确认雅化单位利润}

\begin{exhibitbox}[三家同业的 2026Q1 信号]
\begin{tabularx}{\exhibitboxwidth}{L{2.1cm}R{2.6cm}R{2.8cm}X}
\toprule
公司 & 2026Q1 收入 & 2026Q1 归母净利 & 可确认的直接原因与可比性限制 \\
\midrule
天齐锂业 & 51.28 亿元，+98.44\% & 18.76 亿元，+1,699.12\% & 锂产品均价上升；同时含 SQM 投资收益，不可逐元对比 \\
盛新锂能 & 32.84 亿元，+378.58\% & 4.64 亿元，扭亏 & 锂盐价格上行、印尼出货与量价齐升并存 \\
赣锋锂业 & 未披露于本预告 & 16--21 亿元预告 & 锂价、资源产能释放；另含约 2.59 亿元 PLS 公允价值收益 \\
\bottomrule
\end{tabularx}
\end{exhibitbox}
\sourcenote{天齐锂业《2026 年第一季度报告》（2026-04-28）；盛新锂能《2026 年第一季度报告》（2026-04-30）；赣锋锂业《2026 年第一季度业绩预告》（2026-04-17）。}

三家同业都把锂价上升列为改善因素，支持“行业 $\beta$ 存在”的判断；但天齐的 SQM、赣锋的 PLS 公允价值收益和盛新的印尼出货又说明利润不可横向映射。雅化的矿--产--销协同和成本优化在 H1 预告中是公司表述，不宜在正式 H1 数据缺失时直接认定为已形成可持续成本曲线。

\section{民爆：历史底座存在，H1 景气仍需要公司证据}

工信部关于民爆行业的通知反映安全、数字化与合规是行业约束；锡林郭勒盟 2025 年工业炸药产量同比 +8.61\% 只是矿业地区的活动样本。二者均不能证明雅化 2026 年订单、份额或利润。因此民爆对本事件研究的作用是：等待 H1 分部收入、产品/服务毛利和项目回款来判断其是否构成锂周期的下行保护；在此之前，不给予增长信用。
